\centeredsection{\uppercase{Аннотация}}
В данной курсовой работе проведено исследование процессов прохождения прямоугольных импульсов через фильтр нижних частот (ФНЧ) третьего порядка.

Рассчитана операторная передаточная функция цепи, определены ее нули и полюсы, построены амплитудно- и фазочастотная характеристики.

Выполнено составление уравнений состояния, на основе которых произведен расчет переходной и импульсной характеристик аналитическим методом и численным методом Эйлера.

С использованием спектрального метода оценено влияние полосы пропускания цепи на искажение формы сигналов различной длительности.

Произведено сравнение точного аналитического расчета реакции с приближенным расчетом по спектру.

\vspace{2cm}
\centeredsection{\uppercase{Summary}}
This course work investigates the processes of rectangular pulses passing through a third-order low-pass filter (LPF).

The operational transfer function of the circuit is calculated, its zeros and poles are identified, and the amplitude-frequency and phase-frequency characteristics are plotted.

State equations are formulated, based on which the step and impulse responses are calculated using both the analytical method and Euler's numerical method.

Using the spectral method, the influence of the circuit's bandwidth on the shape distortion of signals with various durations is evaluated.

A comparison is made between the exact analytical calculation of the response and the approximate calculation based on the spectrum.